\section{Modeling}\label{sec:modeling}

After discussing \emph{what a distributed system is}, we now move on to \textbf{how we can structure, describe, and reason about one} (i.e., how we can model it). We divide this section into four major parts:
\begin{itemize}
    \item \textbf{Software architecture}: \emph{what infrastructure supports applications?}
    \item \textbf{Run-time architecture}: \emph{how are components organized and connected?}
    \item \textbf{Interaction model}: \emph{what assumptions can we make about time and communication?}
    \item \textbf{Failure model}: \emph{what can go wrong, and what can we guarantee despite failures?}
\end{itemize}

\newpage

\subsection{Software Architecture and Middleware}\label{sec:modeling--software-architecture-and-middleware}

Before examining specific distributed architectures such as client-server, peer-to-peer, or REST, we need to understand a more fundamental question: \emph{\textbf{how can applications running on different computers communicate and cooperate without having to manage every detail of the underlying platforms?}}

\highspace
Consider two machines:
\begin{itemize}
    \item Machine A runs Linux and an application written in Python.
    \item Machine B runs Windows and an application written in Java.
\end{itemize}
They have different operating systems, programming environments, and potentially different hardware architectures.
Nevertheless, we want their applications to \hl{communicate and cooperate as part of a single distributed system.}
There are \textbf{two main approaches} to organizing the software infrastructure supporting this interaction: \textbf{Network Operating Systems (NOS)} and \textbf{Middleware}. Let's examine them in more detail.

\highspace
\begin{flushleft}
    \textcolor{Green3}{\faIcon{network-wired} \textbf{Network OS-Based Architecture}}
\end{flushleft}
A \definition{Network Operating System (NOS)} is an \textbf{operating system} that \hl{provides facilities for computers to communicate over a network.} In a network OS-based architecture, applications rely directly on networking services provided by their respective operating systems.

\highspace
For example, a programmer may use TCP sockets to exchange messages between two machines. The operating system provides basic facilities such as:
\begin{itemize}
    \item Establishing network connections.
    \item Sending and receiving bytes.
    \item Managing transport protocols such as TCP and UDP.
\end{itemize}
\textcolor{Red2}{\faIcon{exclamation-triangle}} However, these facilities do \hl{not automatically define how the applications should interpret the exchanged data or coordinate their behavior.}

\highspace
Let's imagine another scenario. \emph{Machine A} wants to invoke an operation on \emph{Machine B}:
\begin{lstlisting}[language=python]
result = calculate_price(product_id)
\end{lstlisting}
If \texttt{calculate\_price()} is implemented on \emph{Machine B}, the programmer cannot simply invoke it as an ordinary local function. Instead, the application must implement the required communication logic:
\begin{enumerate}
    \item Establish a connection with \emph{Machine B}.
    \item Serialize the request into a transferable representation.
    \item Send the request across the network.
    \item Wait for and receive a response.
    \item Deserialize the response.
    \item Handle network errors, timeouts, and other failures.
\end{enumerate}
Moreover, the implementations may need adaptations for different programming environments and platform-specific interfaces.

\highspace
Notice that \hl{some of these problems are common to almost every distributed application.} \emph{\textbf{Why should every application developer have to implement them again?}} This motivates the second approach: \textbf{middleware-based architecture}.

\highspace
\begin{figure}[htbp]
    \centering
    \tikzsetnextfilename{network-os-architecture}
    \begin{tikzpicture}[
            font=\small,
            >={Stealth[length=2mm]},
            layer/.style={draw=black!70, thick, rounded corners=2pt,
            minimum width=4.1cm, align=center, inner sep=3pt},
            app/.style={layer, fill=blue!8, minimum height=3.05cm},
            api/.style={layer, fill=orange!15, minimum height=0.55cm},
            nos/.style={layer, fill=green!10, minimum height=1.45cm},
            hw/.style={layer, fill=gray!15, minimum height=0.55cm},
            proto/.style={draw=green!40!black, fill=white, rounded corners=1pt,
            minimum width=1.2cm, minimum height=0.45cm, font=\footnotesize},
            steps/.style={draw=red!60!black, dashed, rounded corners=2pt, fill=white,
            text width=3.6cm, align=left, font=\scriptsize, inner sep=3pt},
            machine/.style={draw=black!40, dashed, rounded corners=4pt, inner sep=4pt},
            pro/.style={text width=5.25cm, align=left, inner sep=4pt,
            font=\scriptsize, execute at begin node={\hyphenpenalty=10000}},
            con/.style={pro},
            idea/.style={draw=blue!50!black, very thick, fill=blue!6, rounded corners=3pt,
                text width=10.4cm, align=center, inner sep=4pt,
            font=\footnotesize, execute at begin node={\hyphenpenalty=10000}},
        ]

        % Machine A: Linux + Python
        \node[hw]  (hwA)  at (2.15,0.30) {Hardware + NIC (e.g.\ x86-64)};
        \node[nos, above=0.12cm of hwA]  (osA) {};
        \node[anchor=north, font=\small\bfseries] at (osA.north) {Network OS: Linux};
        \node[proto] (tcpA) at ($(osA.center)+(-0.7,-0.15)$) {TCP};
        \node[proto] (udpA) at ($(osA.center)+( 0.7,-0.15)$) {UDP};
        \node[font=\scriptsize, anchor=south] at (osA.south) {IP, drivers};
        \node[api, above=0.12cm of osA] (apiA) {Socket API (POSIX sockets)};
        \node[app, above=0.12cm of apiA] (appA) {};
        \node[anchor=north, font=\small\bfseries] at (appA.north) {Python application};
        \node[anchor=north, font=\scriptsize\ttfamily] at ($(appA.north)+(0,-0.42)$)
        {r = calculate\_price(id)};
        \node[steps, anchor=south] (stA) at ($(appA.south)+(0,0.08)$) {%
            \textbf{\color{red!60!black}Developer must write:}\\[1pt]
            1.\ connect to B\\
            2.\ serialize request\\
            3.\ send bytes\\
            4.\ wait \& receive reply\\
            5.\ deserialize result\\
        6.\ handle errors, timeouts};
        \node[machine, fit=(hwA)(appA), label={[font=\bfseries]above:Machine A}] (mA) {};

        % Machine B: Windows + Java
        \node[hw]  (hwB)  at (9.05,0.30) {Hardware + NIC (e.g.\ ARM64)};
        \node[nos, above=0.12cm of hwB]  (osB) {};
        \node[anchor=north, font=\small\bfseries] at (osB.north) {Network OS: Windows};
        \node[proto] (tcpB) at ($(osB.center)+(-0.7,-0.15)$) {TCP};
        \node[proto] (udpB) at ($(osB.center)+( 0.7,-0.15)$) {UDP};
        \node[font=\scriptsize, anchor=south] at (osB.south) {IP, drivers};
        \node[api, above=0.12cm of osB] (apiB) {Socket API (Winsock)};
        \node[app, above=0.12cm of apiB] (appB) {};
        \node[anchor=north, font=\small\bfseries] at (appB.north) {Java application};
        \node[anchor=north, font=\scriptsize\ttfamily] at ($(appB.north)+(0,-0.42)$)
        {double calculatePrice(int)};
        \node[steps, anchor=south] (stB) at ($(appB.south)+(0,0.08)$) {%
            \textbf{\color{red!60!black}Developer must write:}\\[1pt]
            1.\ accept connection\\
            2.\ receive bytes\\
            3.\ deserialize request\\
            4.\ call calculatePrice\\
            5.\ serialize \& send result\\
        6.\ handle errors, timeouts};
        \node[machine, fit=(hwB)(appB), label={[font=\bfseries]above:Machine B}] (mB) {};

        % Network between the two machines
        \node[cloud, cloud puffs=10, aspect=1.8, draw=black!60, fill=cyan!8,
        minimum width=1.7cm, minimum height=1.1cm, font=\footnotesize]
        (net) at ($(osA.east)!0.5!(osB.west)$) {Network};
        \draw[<->, very thick, black!70] (osA.east) -- (net.west);
        \draw[<->, very thick, black!70] (net.east) -- (osB.west);
        \node[font=\scriptsize, align=center, below=0.05cm of net]
        {raw bytes\\(stream or\\datagrams)};

        % What the programmer would like vs. what actually happens
        \draw[->, thick, dashed, red!60!black]
        ($(appA.east)+(0,0.9)$) -- ($(appB.west)+(0,0.9)$)
        node[midway, above, font=\scriptsize, align=center]
        {remote call?\\not a local\\function call}
        node[midway, font=\Large, red!70!black] {$\times$};

        % The problem: the same logic is re-implemented on both sides
        \draw[<->, thick, red!60!black] ($(stA.east)+(0,-0.1)$) -- ($(stB.west)+(0,-0.1)$)
        node[midway, below=1pt, font=\scriptsize\bfseries, align=center, text=red!60!black]
        {same logic,\\written twice,\\in two\\languages};

        % Strengths vs. weaknesses, and the way out
        \node[pro, anchor=north west] (pro) at ($(mA.south west)+(0,-0.3)$) {%
            \textbf{\color{green!40!black}\checkmark\ Strengths}\par\smallskip
            \begin{tabular}{@{}l@{\,}p{4.85cm}@{}}
                $+$ & \raggedright simple: only OS and network, no extra layer \tabularnewline
                $+$ & \raggedright full control of protocol, data format and performance \tabularnewline
                $+$ & \raggedright no dependency on extra infrastructure \tabularnewline
                $+$ & \raggedright low overhead: data goes straight to the OS sockets \tabularnewline
        \end{tabular}};
        \node[con, anchor=north east] (con) at ($(mB.south east)+(0,-0.3)$) {%
            \textbf{\color{red!60!black}$\times$\ Weaknesses}\par\smallskip
            \begin{tabular}{@{}l@{\,}p{4.85cm}@{}}
                $-$ & \raggedright every app re-implements connect, serialize, send, receive, errors \tabularnewline
                $-$ & \raggedright platform differences (POSIX vs.\ Winsock, data formats) left to developers \tabularnewline
                $-$ & \raggedright low-level and error-prone; coordination built ad hoc \tabularnewline
        \end{tabular}};
        \begin{scope}[on background layer]
            \path let \p1 = (pro.south), \p2 = (con.south) in
            coordinate (low) at (0, {min(\y1,\y2)});
            \draw[draw=green!40!black, fill=green!5, rounded corners=2pt]
            (pro.north west) rectangle (pro.east |- low);
            \draw[draw=red!60!black, fill=red!4, rounded corners=2pt]
            (con.north west) rectangle (con.east |- low);
        \end{scope}
        \path let \p1 = (pro.south), \p2 = (con.south) in
        node[idea, anchor=north] (idea) at (5.6, {min(\y1,\y2)-0.5cm}) {%
            \textbf{Why rewrite it in every application?}
            Move the common communication code into a reusable layer
        between applications and OS $\Rightarrow$ \textbf{Middleware}};
        \draw[->, very thick, red!60!black] (con.south |- low) -- (con.south |- idea.north);
    \end{tikzpicture}
    \caption{\textbf{Network OS-based architecture.} Each application relies directly on the
        networking facilities (TCP/UDP sockets) of its own operating system. The OS
        only moves bytes: connecting, serializing, sending, receiving, deserializing
        and handling failures must be implemented by each application, on each
        platform. This gives full control, but the same generic logic is rewritten
    again and again, which motivates middleware.}
    \label{fig:network-os-architecture}
\end{figure}

\highspace
\begin{flushleft}
    \textcolor{Green3}{\faIcon{layer-group} \textbf{Middleware-Based Architecture}}
\end{flushleft}
\definition{Middleware} is a \textbf{software layer} positioned logically \hl{between distributed applications and the underlying operating-system/network facilities.} Its purpose is to \hl{provide reusable services and abstractions that simplify the development of distributed applications.}

\highspace
\textcolor{Green3}{\faIcon[regular]{lightbulb}} The central idea is that applications communicate through a \textbf{common, well-defined API} instead of dealing directly with all the underlying networking details.

\highspace
For example, consider a middleware supporting Remote Procedure Calls (RPC). An application may invoke:
\begin{lstlisting}[language=python]
result = remote_service.calculate_price(product_id)
\end{lstlisting}
The middleware can take responsibility for: (1) serializing arguments, (2) transmitting requests, (3) receiving responses, and (4) reconstructing results. The \hl{programmer works with a higher-level abstraction}, like invoking a remote operation rather than manually transmitting bytes over a socket.

\highspace
The middleware also \textbf{hides many differences between platforms} by providing compatible interfaces and communication protocols.
Note: the operating systems are still involved. Middleware does not replace their networking facilities, \textbf{it builds higher-level abstractions on top of them}.

\highspace
\begin{table}[!htp]
    \centering
    \begin{tabular}{@{} >{\raggedright\arraybackslash}p{.22\linewidth} >{\raggedright\arraybackslash}p{.34\linewidth} >{\raggedright\arraybackslash}p{.34\linewidth} @{}}
        \toprule
        \textbf{Aspect} & \textbf{Network OS-Based} & \textbf{Middleware-Based} \\
        \midrule
        Communication abstraction & OS networking facilities & High-level APIs \\[1.8em]
        Platform differences & Largely handled by application developers & Partly hidden by middleware \\[1.8em]
        Coordination services & Often developed by applications & Can be provided as reusable services \\[1.8em]
        Development complexity & More low-level concerns & More infrastructure functionality reused \\[1.8em]
        Examples & Direct socket programming & RPC frameworks, message brokers \\
        \bottomrule
    \end{tabular}
\end{table}

\highspace
Middleware simplifies many aspects of distribution, but it does \hl{not make remote communication identical to local execution.} Network latency, partial failures, and inconsistent system states remain relevant.

\begin{figure}[htbp]
    \centering
    \tikzsetnextfilename{middleware-architecture}
    \begin{tikzpicture}[
            font=\small,
            >={Stealth[length=2mm]},
            layer/.style={draw=black!70, thick, rounded corners=2pt,
            minimum width=4.1cm, align=center, inner sep=3pt},
            app/.style={layer, fill=blue!8, minimum height=1.55cm},
            mwapi/.style={layer, fill=violet!20, minimum height=0.5cm, font=\footnotesize},
            mw/.style={layer, fill=violet!8, draw=violet!60!black, very thick,
            minimum height=2.2cm},
            api/.style={layer, fill=orange!15, minimum height=0.5cm, font=\footnotesize},
            nos/.style={layer, fill=green!10, minimum height=1.25cm},
            hw/.style={layer, fill=gray!15, minimum height=0.5cm, font=\footnotesize},
            proto/.style={draw=green!40!black, fill=white, rounded corners=1pt,
            minimum width=1.2cm, minimum height=0.4cm, font=\footnotesize},
            svc/.style={draw=violet!60!black, fill=white, rounded corners=1pt,
            minimum width=1.85cm, minimum height=0.4cm, font=\scriptsize},
            machine/.style={draw=black!40, dashed, rounded corners=4pt, inner sep=4pt},
            pro/.style={text width=5.25cm, align=left, inner sep=4pt,
            font=\scriptsize, execute at begin node={\hyphenpenalty=10000}},
            idea/.style={draw=blue!50!black, very thick, fill=blue!6, rounded corners=3pt,
                text width=10.4cm, align=center, inner sep=4pt,
            font=\footnotesize, execute at begin node={\hyphenpenalty=10000}},
        ]

        % Machine A: Linux + Python
        \node[hw]  (hwA)  at (2.15,0.25) {Hardware + NIC (e.g.\ x86-64)};
        \node[nos, above=0.12cm of hwA]  (osA) {};
        \node[anchor=north, font=\small\bfseries] at (osA.north) {Network OS: Linux};
        \node[proto] (tcpA) at ($(osA.center)+(-0.7,-0.12)$) {TCP};
        \node[proto] (udpA) at ($(osA.center)+( 0.7,-0.12)$) {UDP};
        \node[api, above=0.12cm of osA] (apiA) {Socket API (POSIX sockets)};
        \node[mw, above=0.12cm of apiA] (mwA) {};
        \node[anchor=north, font=\small, text=violet!60!black] at (mwA.north)
        {\textbf{Middleware} \footnotesize(e.g.\ gRPC)};
        \node[svc] at ($(mwA.center)+(-0.97, 0.20)$) {(de)serialize};
        \node[svc] at ($(mwA.center)+( 0.97, 0.20)$) {send/receive};
        \node[svc] at ($(mwA.center)+(-0.97,-0.27)$) {failures};
        \node[svc] at ($(mwA.center)+( 0.97,-0.27)$) {naming};
        \node[svc] at ($(mwA.center)+(-0.97,-0.74)$) {security};
        \node[svc] at ($(mwA.center)+( 0.97,-0.74)$) {transactions};
        \node[mwapi, above=0.12cm of mwA] (mapiA) {Middleware API (common)};
        \node[app, above=0.12cm of mapiA] (appA) {};
        \node[anchor=north, font=\small\bfseries] at (appA.north) {Python application};
        \node[anchor=north, font=\scriptsize\ttfamily, align=center]
        at ($(appA.north)+(0,-0.42)$)
        {r = remote\_service.\\\phantom{r = }calculate\_price(id)};
        \node[anchor=south, font=\scriptsize\itshape] at (appA.south) {business logic only};
        \node[machine, fit=(hwA)(appA), label={[font=\bfseries]above:Machine A}] (mA) {};

        % Machine B: Windows + Java
        \node[hw]  (hwB)  at (9.05,0.25) {Hardware + NIC (e.g.\ ARM64)};
        \node[nos, above=0.12cm of hwB]  (osB) {};
        \node[anchor=north, font=\small\bfseries] at (osB.north) {Network OS: Windows};
        \node[proto] (tcpB) at ($(osB.center)+(-0.7,-0.12)$) {TCP};
        \node[proto] (udpB) at ($(osB.center)+( 0.7,-0.12)$) {UDP};
        \node[api, above=0.12cm of osB] (apiB) {Socket API (Winsock)};
        \node[mw, above=0.12cm of apiB] (mwB) {};
        \node[anchor=north, font=\small, text=violet!60!black] at (mwB.north)
        {\textbf{Middleware} \footnotesize(e.g.\ gRPC)};
        \node[svc] at ($(mwB.center)+(-0.97, 0.20)$) {(de)serialize};
        \node[svc] at ($(mwB.center)+( 0.97, 0.20)$) {send/receive};
        \node[svc] at ($(mwB.center)+(-0.97,-0.27)$) {failures};
        \node[svc] at ($(mwB.center)+( 0.97,-0.27)$) {naming};
        \node[svc] at ($(mwB.center)+(-0.97,-0.74)$) {security};
        \node[svc] at ($(mwB.center)+( 0.97,-0.74)$) {transactions};
        \node[mwapi, above=0.12cm of mwB] (mapiB) {Middleware API (common)};
        \node[app, above=0.12cm of mapiB] (appB) {};
        \node[anchor=north, font=\small\bfseries] at (appB.north) {Java application};
        \node[anchor=north, font=\scriptsize\ttfamily] at ($(appB.north)+(0,-0.42)$)
        {double calculatePrice(int)};
        \node[anchor=south, font=\scriptsize\itshape] at (appB.south) {business logic only};
        \node[machine, fit=(hwB)(appB), label={[font=\bfseries]above:Machine B}] (mB) {};

        % 1. Application level: a remote call that looks like a method call
        \draw[->, thick, dashed, green!40!black] (appA.east) -- (appB.west)
        node[midway, above, font=\scriptsize, align=center]
        {remote call\\through the\\middleware}
        node[midway, below, font=\Large, green!40!black] {\checkmark};

        % 2. Middleware level: the two components speak a compatible protocol
        \draw[<->, very thick, dashed, violet!60!black] (mwA.east) -- (mwB.west)
        node[midway, above, font=\scriptsize, align=center, text=violet!60!black]
        {compatible\\protocol}
        node[midway, below, font=\scriptsize, align=center, text=violet!60!black]
        {(e.g.\ HTTP/2\\+ Protobuf)};

        % 3. OS level: bytes still travel through the OS networking facilities
        \node[cloud, cloud puffs=10, aspect=1.8, draw=black!60, fill=cyan!8,
        minimum width=1.7cm, minimum height=1.0cm, font=\footnotesize]
        (net) at ($(osA.east)!0.5!(osB.west)$) {Network};
        \draw[<->, very thick, black!70] (osA.east) -- (net.west);
        \draw[<->, very thick, black!70] (net.east) -- (osB.west);
        \node[font=\scriptsize, align=center, below=0.05cm of net]
        {OS still moves\\the bytes};

        % Benefits vs. what remains, and the take-away
        \node[pro, anchor=north west] (pro) at ($(mA.south west)+(0,-0.3)$) {%
            \textbf{\color{green!40!black}\checkmark\ Benefits}\par\smallskip
            \begin{tabular}{@{}l@{\,}p{4.85cm}@{}}
                $+$ & \raggedright communication code written once and reused by every app \tabularnewline
                $+$ & \raggedright higher-level abstraction: call an operation, not send bytes \tabularnewline
                $+$ & \raggedright platform differences partly hidden by common API and protocol \tabularnewline
                $+$ & \raggedright reusable services: naming, security, transactions, messaging \tabularnewline
        \end{tabular}};
        \node[pro, anchor=north east] (con) at ($(mB.south east)+(0,-0.3)$) {%
            \textbf{\color{orange!70!black}!\ Still to consider}\par\smallskip
            \begin{tabular}{@{}l@{\,}p{4.85cm}@{}}
                $-$ & \raggedright network latency: a remote call is slower than a local one \tabularnewline
                $-$ & \raggedright partial failures: the other side may crash or be unreachable \tabularnewline
                $-$ & \raggedright inconsistent states between the components \tabularnewline
                $-$ & \raggedright an extra layer to deploy, configure and depend on \tabularnewline
        \end{tabular}};
        \begin{scope}[on background layer]
            \path let \p1 = (pro.south), \p2 = (con.south) in
            coordinate (low) at (0, {min(\y1,\y2)});
            \draw[draw=green!40!black, fill=green!5, rounded corners=2pt]
            (pro.north west) rectangle (pro.east |- low);
            \draw[draw=orange!70!black, fill=orange!5, rounded corners=2pt]
            (con.north west) rectangle (con.east |- low);
        \end{scope}
        \node[idea, anchor=north] (idea) at ($(5.6,0)+(0,-0.5)+(0,0 |- low)$) {%
            \textbf{Middleware builds on the OS, it does not replace it.}
            It simplifies distributed programming, but a remote call is
        \emph{not} identical to a local one};
        \draw[->, very thick, orange!70!black] (con.south |- low) -- (con.south |- idea.north);
    \end{tikzpicture}
    \caption{\textbf{Middleware-based architecture.} A middleware layer sits between each
        application and its operating system. The applications use a common,
        high-level API (e.g.\ a remote procedure call), while the two middleware
        components take care of serialization, transmission, failures and other
        reusable services, talking to each other through a compatible protocol.
    The operating systems still move the bytes over the network.}
    \label{fig:middleware-architecture}
\end{figure}

\newpage

\begin{flushleft}
    \textcolor{Green3}{\faIcon{layer-group} \textbf{Middleware as a Functional Abstraction}}
\end{flushleft}
We have described \textbf{where middleware fits} in a distributed system, and now we need to understand \textbf{what services it provides}.

\highspace
A fundamental property is that middleware typically provides \textbf{business-\break unaware services}. This means that middleware does not need to understand the specific business logic of an application. Consider an e-commerce platform:
\begin{itemize}
    \item The application knows what a customer, product, order, or payment represents.
    \item The middleware does not need to know these concepts.
\end{itemize}
Instead, it \hl{provides generic mechanisms for communicating, coordinating, and managing the distributed components} implementing that business logic.
We can \textbf{classify middleware functionality} into three main categories:
\begin{enumerate}
    \item \important{Communication and Coordination Services.} These services determine \hl{how distributed components exchange information and coordinate their activities.} Communication can be:
        \begin{itemize}
            \item \textbf{Synchronous}, the \textbf{caller waits} for a response before continuing.
            \item \textbf{Asynchronous}, the \textbf{sender} can \textbf{continue} without waiting for a response.
        \end{itemize}
        Also, communication can be:
        \begin{itemize}
            \item \textbf{Point-to-point}, where a message is sent to a \textbf{specific recipient}.
            \item \textbf{Multicast}, where a message is sent to \textbf{multiple recipients} simultaneously.
        \end{itemize}
        In addition to communication, middleware can also provide \textbf{coordination primitives}, such as distributed locks, synchronization mechanisms, or coordination protocols.

        \begin{figure}[H]
            \centering
            \tikzsetnextfilename{middleware-communication-coordination}
            \begin{tikzpicture}[
                    font=\small,
                    >={Stealth[length=2mm]},
                    panel/.style={draw=violet!60!black, thick, rounded corners=3pt, fill=violet!4},
                    title/.style={anchor=north west, font=\small\bfseries, text=violet!60!black},
                    comp/.style={draw=black!70, thick, rounded corners=2pt, fill=blue!8,
                    minimum width=1.5cm, minimum height=0.5cm, font=\footnotesize, align=center},
                    mwnode/.style={comp, fill=violet!15, draw=violet!60!black},
                    life/.style={draw=black!40, dashed},
                    msg/.style={->, thick},
                    note/.style={font=\scriptsize, align=center},
                ]

                % Synchronous
                \draw[panel] (0,0) rectangle (5.55,3.9);
                \node[title] at (0.1,3.85) {Synchronous {\normalfont\footnotesize(e.g.\ blocking RPC)}};
                \node[comp] (cal) at (1.2,3.05) {Caller};
                \node[comp] (cle) at (4.3,3.05) {Callee};
                \draw[life] (cal.south) -- (1.2,0.2);
                \draw[life] (cle.south) -- (4.3,0.2);
                \fill[red!30, draw=red!60!black] (1.12,2.55) rectangle (1.28,0.85);
                \fill[blue!20, draw=blue!50!black] (4.22,2.25) rectangle (4.38,1.2);
                \draw[msg] (1.28,2.55) -- (4.22,2.25) node[midway, above, note] {request};
                \draw[msg] (4.22,1.2) -- (1.28,0.9) node[midway, below, note] {response};
                \node[note, anchor=west, text=red!60!black] at (1.35,1.75) {caller is blocked:\\it waits for the reply};

                % Asynchronous
                \begin{scope}[xshift=5.85cm]
                    \draw[panel] (0,0) rectangle (5.55,3.9);
                    \node[title] at (0.1,3.85) {Asynchronous {\normalfont\footnotesize(e.g.\ message queue)}};
                    \node[comp] (snd) at (1.0,3.05) {Sender};
                    \node[comp] (rcv) at (4.5,3.05) {Receiver};
                    \draw[life] (snd.south) -- (1.0,0.2);
                    \draw[life] (rcv.south) -- (4.5,0.2);
                    \foreach \i in {0,1,2}
                    \draw[draw=violet!60!black, fill=violet!15] (2.4+0.25*\i,2.3) rectangle ++(0.22,0.35);
                    \node[note, text=violet!60!black, above] at (2.735,2.65) {queue};
                    \fill[green!30, draw=green!40!black] (0.92,2.5) rectangle (1.08,0.4);
                    \fill[blue!20, draw=blue!50!black] (4.42,1.6) rectangle (4.58,0.75);
                    \draw[msg] (1.08,2.5) -- (2.38,2.45) node[midway, below, note] {send};
                    \draw[msg, dashed] (3.17,2.4) -- (4.42,1.6) node[midway, above right=-2pt, note] {later};
                    \node[note, anchor=west, text=green!40!black] at (1.15,1.3) {sender continues\\without waiting};
                \end{scope}

                % Point-to-point
                \draw[panel] (0,-3.4) rectangle (3.6,-0.3);
                \node[title] at (0.1,-0.35) {Point-to-point};
                \node[comp] (ord) at (1.8,-1.3) {Order service};
                \node[comp] (pay) at (1.8,-2.75) {Payment service};
                \draw[msg] (ord) -- (pay) node[midway, right, note] {pay(order)};
                \node[note, text=black!60] at (1.8,-3.18) {one specific recipient};

                % Multicast
                \draw[panel] (3.9,-3.4) rectangle (7.5,-0.3);
                \node[title] at (4.0,-0.35) {Multicast};
                \node[comp] (cfg) at (5.7,-1.3) {Config service};
                \foreach \x/\n in {4.55/1, 5.7/2, 6.85/3}
                \node[comp, minimum width=0.9cm] (s\n) at (\x,-2.75) {Svc \n};
                \foreach \n in {1,2,3}
                \draw[msg] (cfg.south) -- (s\n.north);
                \node[note, text=black!60] at (5.7,-3.18) {``config updated'' to many};

                % Coordination: distributed lock
                \draw[panel] (7.8,-3.4) rectangle (11.4,-0.3);
                \node[title] at (7.9,-0.35) {Coordination};
                \node[mwnode] (lock) at (9.6,-1.3) {Lock service};
                \node[comp, minimum width=1.2cm] (c1) at (8.6,-2.75) {Comp.\ 1};
                \node[comp, minimum width=1.2cm] (c2) at (10.6,-2.75) {Comp.\ 2};
                \draw[msg, green!40!black] (c1.north) -- (lock.south west)
                node[pos=0.35, right, note, text=green!40!black] {granted\,\checkmark};
                \draw[msg, red!60!black, dashed] (c2.north) -- (lock.south east)
                node[pos=0.7, left, note, text=red!60!black] {waits};
                \node[note, text=black!60] at (9.6,-3.18) {one holder at a time};
            \end{tikzpicture}
            % \caption{Communication and coordination services offered by middleware.
            %     Top: in a synchronous interaction the caller waits for the response,
            %     while in an asynchronous one the sender continues and the message is
            %     delivered later. Bottom: a message can go to one specific recipient
            %     (point-to-point) or to many (multicast); coordination primitives such as
            % a distributed lock let components agree on who may act.}
            \label{fig:middleware-communication-coordination}
        \end{figure}

    \item \important{Special Application Services.} These are \hl{reusable, higher-level functionalities} provided by middleware to \hl{solve common application problems.}

        \begin{examplebox}[: E-commerce Middleware Services]
            For example, in an e-commerce application, the middleware might provide:
            \begin{itemize}
                \item \textbf{Distributed transaction management}: coordinate operations across several components so that they respect transaction guarantees.
                \item \textbf{Messaging services}: store, route, and deliver messages between distributed components.
                \item \textbf{Notification services}: inform interested components when relevant events occur.
                \item \textbf{Workflow and groupware services}: support coordinated activities involving multiple participants or applications.
            \end{itemize}
            For instance, consider an order that requires updating an inventory service and a payment service.
            The middleware might provide distributed transaction mechanisms to coordinate the updates.
            However, it does not determine whether a customer is allowed to purchase a particular product. That decision belongs to the application.
            The distinction is:
            \begin{equation*}
                \underbrace{\text{How to coordinate updates}}_{\text{Middleware}}
                \neq
                \underbrace{\text{Which updates are valid}}_{\text{Application logic}}
            \end{equation*}
            So \textbf{middleware} provides the \textbf{mechanisms} (e.g., message queues, transaction coordination), while the \textbf{application} decides \textbf{how to use them} (e.g., send an order confirmation after payment is successful).
            \begin{center}
                \tikzsetnextfilename{middleware-application-services}
                \tikzwidth{\textwidth}
                \begin{tikzpicture}[
                        font=\small,
                        >={Stealth[length=2mm]},
                        comp/.style={draw=black!70, thick, rounded corners=2pt, fill=blue!8,
                        minimum width=2.4cm, minimum height=0.8cm, font=\footnotesize, align=center},
                        chip/.style={draw=violet!60!black, fill=white, rounded corners=2pt,
                        text width=2.35cm, minimum height=1.0cm, align=center, font=\scriptsize},
                        hot/.style={chip, very thick, fill=violet!18},
                        role/.style={font=\scriptsize\itshape, align=right, anchor=north east},
                        msg/.style={->, thick},
                        note/.style={font=\scriptsize, align=center},
                    ]

                    % Application layer: knows the business
                    \draw[thick, draw=black!70, fill=blue!8, rounded corners=3pt] (0,3.65) rectangle (11.4,5.05);
                    \node[anchor=north west, font=\small\bfseries] at (0.1,5.0) {Order application};
                    \node[role, text=blue!50!black] at (11.3,5.0)
                    {business-aware:\\decides \textbf{which} updates are valid};
                    \node[anchor=west, font=\footnotesize] at (0.2,4.0)
                    {Business rule: \emph{may this customer buy this product?}};

                    % Middleware layer: generic, reusable services
                    \draw[very thick, draw=violet!60!black, fill=violet!6, rounded corners=3pt]
                    (0,0.9) rectangle (11.4,2.85);
                    \node[anchor=north west, font=\small\bfseries, text=violet!60!black] at (0.1,2.8)
                    {Middleware};
                    \node[role, text=violet!60!black] at (11.3,2.8)
                    {business-unaware:\\decides \textbf{how} to coordinate updates};
                    \node[chip] (msgq) at (1.45,1.5) {\textbf{Messaging}\\store, route, deliver};
                    \node[hot]  (tx)   at (4.28,1.5) {\textbf{Distributed\\transactions}};
                    \node[chip] (ntf)  at (7.11,1.5) {\textbf{Notifications}\\inform on events};
                    \node[chip] (wf)   at (9.94,1.5) {\textbf{Workflow /}\\\textbf{groupware}};

                    \draw[msg] (4.28,3.65) -- (tx.north)
                    node[pos=0.22, right, note] {place order (update stock + charge)};

                    % Services touched by the transaction
                    \node[comp] (inv) at (2.6,-0.7)  {Inventory service};
                    \node[comp] (pay) at (5.96,-0.7) {Payment service};
                    \draw[msg] (tx.south) -- (inv.north) node[midway, above left=-1pt, note] {reserve stock};
                    \draw[msg] (tx.south) -- (pay.north) node[midway, above right=-1pt, note] {charge};
                    \node[draw=violet!60!black, dashed, rounded corners=2pt, fill=violet!4,
                        text width=3.3cm, align=left, font=\scriptsize, anchor=west,
                    execute at begin node={\hyphenpenalty=10000}] at (7.75,-0.7)
                    {\textbf{All or nothing:} both updates commit, or both are rolled back
                    (e.g.\ payment fails $\Rightarrow$ stock is released).};

                    % The key distinction
                    % \node[font=\footnotesize] at (5.7,-2.3)
                    % {$\underbrace{\text{How to coordinate updates}}_{\text{Middleware}}
                    %     \;\neq\;
                    % \underbrace{\text{Which updates are valid}}_{\text{Application logic}}$};
                \end{tikzpicture}
            \end{center}
        \end{examplebox}

        \newpage

    \item \important{Management Services.} These services address the \textbf{operational complexity} of distributed systems. Three fundamental examples are:
        \begin{itemize}
            \item \textbf{Naming}: locating and identifying distributed resources and services.
            \item \textbf{Security}: providing mechanism for authentication, authorization, and protected communication.
            \item \textbf{Failure handling}: detecting or reporting failures and supporting recovery mechanisms.
        \end{itemize}

        \begin{figure}[H]
            \centering
            \tikzsetnextfilename{middleware-management-services}
            \begin{tikzpicture}[
                    font=\small,
                    >={Stealth[length=2mm]},
                    panel/.style={draw=violet!60!black, thick, rounded corners=3pt, fill=violet!4},
                    title/.style={anchor=north west, font=\small\bfseries, text=violet!60!black},
                    comp/.style={draw=black!70, thick, rounded corners=2pt, fill=blue!8,
                    minimum width=1.5cm, minimum height=0.55cm, font=\footnotesize, align=center},
                    mwnode/.style={comp, fill=violet!15, draw=violet!60!black,
                    execute at begin node={\hyphenpenalty=10000}},
                    msg/.style={->, thick},
                    note/.style={font=\scriptsize, align=center},
                ]

                % Naming / discovery
                \draw[panel] (0,0) rectangle (11.4,3.7);
                \node[title] at (0.1,3.65) {Naming / discovery};
                \node[anchor=north east, font=\scriptsize\itshape, text=violet!60!black, align=right]
                at (11.3,3.65) {same logical name even if the service moves\\$\Rightarrow$ location transparency};

                \node[comp] (cli) at (1.3,2.0) {Client\\[1pt]\scriptsize\ttfamily authentication-\\[-2pt]\scriptsize\ttfamily service};
                \node[mwnode, text width=3.0cm] (ns) at (5.75,2.0)
                {\textbf{Naming service}\\[2pt]
                    \scriptsize\ttfamily authentication-service\\
                \scriptsize$\mapsto$ \ttfamily 10.0.1.7:8443};
                \node[comp, dashed, draw=black!40, fill=gray!8, text=black!50] (old) at (9.7,2.3)
                {Machine X\\\scriptsize(old location)};
                \node[font=\Large, red!60!black] at (old.north east) {$\times$};
                \node[comp] (new) at (9.7,0.7) {Machine Y\\\scriptsize auth service\\\scriptsize\ttfamily 10.0.1.7};
                \draw[->, dashed, black!40] (old.south) -- (new.north) node[midway, right, note, text=black!50] {moved};

                \draw[msg] ($(cli.east)+(0,0.3)$) -- ($(ns.west)+(0,0.3)$)
                node[midway, above, note] {1.\ resolve};
                \draw[msg] ($(ns.west)+(0,-0.3)$) -- ($(cli.east)+(0,-0.3)$)
                node[midway, below, note] {2.\ endpoint};
                \draw[msg] (cli.south) |- (new.west)
                node[pos=0.75, above, note] {3.\ connect};
                \draw[->, dashed, violet!60!black] (new.north west) -- (ns.south east)
                node[midway, below left=-2pt, note, text=violet!60!black] {0.\ register};

                % Security: checks on the way to the service, over a protected channel
                \draw[panel] (0,-4.6) rectangle (5.55,-0.3);
                \node[title] at (0.1,-0.35) {Security};
                \node[comp, minimum width=2.71cm] (sc) at (1.6,-1.15) {Client};
                \node[mwnode, text width=2.5cm, font=\scriptsize] (au) at (1.6,-2.1)
                {\textbf{1.\ Authentication}\\who are you?};
                \node[mwnode, text width=2.5cm, font=\scriptsize] (az) at (1.6,-3.15)
                {\textbf{2.\ Authorization}\\may you do this?};
                \node[comp, minimum width=2.71cm] (ss) at (1.6,-4.1) {Service};
                \draw[msg] (sc) -- (au);
                \draw[msg] (au) -- (az);
                \draw[msg] (az) -- (ss);
                % the whole exchange travels over a protected channel (bracket broken by the padlock)
                \draw[very thick, green!40!black] (3.1,-1.15) -- (3.25,-1.15) -- (3.25,-2.2);
                \draw[very thick, green!40!black] (3.25,-2.95) -- (3.25,-4.1) -- (3.1,-4.1);
                \begin{scope}[shift={(3.25,-2.625)}] % padlock icon
                    \draw[thick, fill=white] (-0.15,0.06) arc (180:0:0.15 and 0.17);
                    \draw[thick, fill=yellow!60!orange] (-0.21,-0.17) rectangle (0.21,0.09);
                \end{scope}
                \node[note, align=left, anchor=west, text=green!40!black] at (3.5,-2.625)
                {\textbf{3.\ Protected}\\\textbf{channel}\\(e.g.\ TLS):\\nobody can\\read or alter\\the data};

                % Failure handling: detect, report, fail over
                \draw[panel] (5.85,-4.6) rectangle (11.4,-0.3);
                \node[title] at (5.95,-0.35) {Failure handling};
                \node[comp, minimum width=1.3cm] (fc) at (8.6,-1.15) {Client};
                \node[mwnode] (fd) at (8.6,-2.55) {Failure detector};
                \node[comp, minimum width=1.4cm, draw=red!60!black, fill=red!8] (r1) at (6.95,-4.15)
                {Replica 1 {\color{red!60!black}$\times$}};
                \node[comp, minimum width=1.4cm] (r2) at (10.45,-4.15) {Replica 2};
                % heartbeats
                \draw[->, thick, dashed, black!35] ($(r1.north)+(0.4,0)$) -- (fd.south west)
                node[pos=0.3, right, note, text=red!60!black] {no\\heartbeat};
                \draw[->, thick, dashed, green!40!black] ($(r2.north)+(-0.4,0)$) -- (fd.south east)
                node[pos=0.7, left, note] {heartbeat};
                % report and fail over
                \draw[->, thick, violet!60!black] (fd) -- (fc)
                node[midway, right, note, text=violet!60!black] {report failure};
                \draw[->, thick, dashed, red!60!black] (fc.west) -| (r1.north)
                node[pos=0.75, left, note, align=right] {before};
                \draw[->, very thick, green!40!black] (fc.east) -| (r2.north)
                node[pos=0.75, right, note, align=left] {after\\fail\\over};
            \end{tikzpicture}
        \end{figure}

        \begin{examplebox}[: Naming and Discovery]
            For example, suppose a client needs to communicate with an \textbf{authentication service}.
            Instead of hardcoding the IP address of a particular server, it might use a logical service name:
            \texttt{authentication-service}.

            \highspace
            A naming or discovery mechanism resolves that name to an appropriate endpoint.
            If the service moves to another machine, the client may continue using the same logical name.
            This connects directly to the concept of \textbf{location transparency} (\autopageref{item:location-transparency}).

            \highspace
            Middleware can \textbf{help achieve transparency}, although how much distribution can actually be hidden depends on its design and on the guarantees offered by the system.
        \end{examplebox}
\end{enumerate}

\newpage

\begin{examplebox}
    Some examples of middleware:
    \begin{itemize}
        \item \textbf{gRPC}: remote procedure calls.
        \item \textbf{MPI}: message passing between processes, especially in parallel computing.
        \item \textbf{Apache Kafka}: distributed event streaming and messaging.
        \item \textbf{CORBA}: distributed object interaction.
    \end{itemize}
    These technologies solve different problems, but they share the principle of \hl{providing reusable abstractions above lower-level networking facilities.}

    \highspace
    MPI, for example, is a specialized form of communication middleware. It offers operations such as \texttt{MPI\_Send}, \texttt{MPI\_Recv}, and collective communication primitives instead of requiring developers to build their own network communication protocols.

    \highspace
    Note: middleware does not necessarily need to be a separate server. It might be \hl{implemented as libraries inside application processes}, as external infrastructure such as message brokers, or as a combination of both.
\end{examplebox}

\highspace
\begin{takeawaysbox}
    Network Operating Systems provide basic communication facilities (i.e., moving bytes), while middleware provides higher-level reusable abstractions (i.e., moving objects, messages, or events).

    \highspace
    \textbf{Middleware provides three broad families of services}:
    \begin{enumerate}
        \item \textbf{Communication and Coordination}: exchanging messages and coordinating components (e.g., RPC, message queues, distributed locks).
        \item \textbf{Special application services}: transactions, messaging, notifications, workflows.
        \item \textbf{Management services}: naming, security, failure handling.
    \end{enumerate}
    Its main \textbf{advantage} is \hl{reducing the complexity of developing distributed applications} by providing \hl{common, platform-independent abstractions and reusable services.}

    In general, middleware simplifies the development of distributed applications, but does \textbf{not eliminate the fundamental challenges of distribution}.

    \begin{center}
        \tikzsetnextfilename{middleware-key-takeaways}
        \tikzwidth{\textwidth}
        \begin{tikzpicture}[
                font=\small,
                >={Stealth[length=2mm]},
                box/.style={draw=#1, thick, rounded corners=3pt, fill=#1!5},
                title/.style={anchor=north west, font=\small\bfseries},
                lay/.style={draw=black!60, rounded corners=1pt, minimum width=2.9cm,
                minimum height=0.42cm, font=\scriptsize, inner sep=1pt},
                txt/.style={anchor=north west, font=\scriptsize, text width=#1, align=left, inner sep=1pt,
                execute at begin node={\hyphenpenalty=10000}},
                chip/.style={draw=violet!60!black, fill=white, rounded corners=2pt, text width=3.25cm,
                    minimum height=1.25cm, align=center, font=\scriptsize,
                execute at begin node={\hyphenpenalty=10000}},
            ]

            % 1. Network OS-based vs. middleware-based
            % --- Network OS-based
            \draw[box=green!40!black] (0,3.95) rectangle (5.1,7.35);
            \node[title, text=green!40!black] at (0.1,7.3) {Network OS-based};
            \node[anchor=north west, font=\scriptsize\itshape] at (0.1,6.9)
            {$\Rightarrow$ basic communication facilities};
            \node[lay, fill=blue!8]   (a1) at (2.55,5.97) {App + communication code};
            \node[lay, fill=green!12, below=1pt of a1] (a2) {Network OS};
            \node[txt=4.85cm] at (0.12,4.95)
            {$-$ raw sockets: bytes over TCP/UDP\\
                $-$ same logic rewritten by every app\\
            $-$ platform differences left to the app};

            % --- Middleware-based
            \draw[box=violet!60!black] (6.3,3.95) rectangle (11.4,7.35);
            \node[title, text=violet!60!black] at (6.4,7.3) {Middleware-based};
            \node[anchor=north west, font=\scriptsize\itshape] at (6.4,6.9)
            {$\Rightarrow$ higher-level reusable abstractions};
            \node[lay, fill=blue!8]    (b1) at (8.85,6.2) {App (business logic only)};
            \node[lay, fill=violet!15, draw=violet!60!black, below=1pt of b1] (b2) {Middleware};
            \node[lay, fill=green!12,  below=1pt of b2] (b3) {Network OS};
            \node[txt=4.85cm] at (6.42,4.95)
            {$+$ common, well-defined API\\
                $+$ reusable services, written once\\
            $+$ platform differences partly hidden};

            \draw[->, very thick, violet!60!black] (5.15,5.55) -- (6.25,5.55)
            node[midway, above, font=\scriptsize, align=center] {factor\\out}
            node[midway, below, font=\scriptsize, align=center] {common\\code};

            % 2. The three families of middleware services
            \draw[box=violet!60!black] (0,0.5) rectangle (11.4,3.65);
            \node[title, text=violet!60!black] at (0.1,3.6) {Three families of middleware services};
            \node[anchor=north east, font=\scriptsize\itshape, text=violet!60!black] at (11.3,3.6)
            {business-unaware: \textbf{how}, not \textbf{which}};
            \node[chip] at (1.9,2.4)
            {\textbf{1.\ Communication \&}\\\textbf{coordination}\\sync/async, locks,\\
            point-to-point/multicast};
            \node[chip] at (5.7,2.4)
            {\textbf{2.\ Special application}\\\textbf{services}\\transactions, messaging,
            notifications, workflows};
            \node[chip] at (9.5,2.4)
            {\textbf{3.\ Management}\\\textbf{services}\\naming, security,\\failure handling};
            \node[font=\scriptsize, align=center] at (5.7,1.05)
            {\textbf{Examples:} gRPC (RPC) $\cdot$ MPI (message passing) $\cdot$\\
                Kafka (event streaming) $\cdot$ CORBA (distributed objects)\\
            as libraries inside the processes, as external brokers, or both};

            % 3. What to keep in mind
            \draw[box=orange!70!black] (0,-1.45) rectangle (11.4,0.2);
            \node[title, text=orange!70!black] at (0.1,0.15) {! Not a silver bullet};
            \node[txt=11.0cm] at (0.15,-0.38)
            {Middleware simplifies the development of distributed applications, but it does
                \emph{not} remove the challenges of distribution: network latency, partial
            failures and inconsistent states remain.};
        \end{tikzpicture}
    \end{center}
    In summary, a network OS offers only basic communication facilities, while middleware factors the common logic into higher-level reusable abstractions. Its services fall into three business-unaware families. Middleware simplifies distribution without eliminating its challenges.
\end{takeawaysbox}